\subsection{默认配置的有限动作与虚拟时间上界}
\label{sec:default-finite-bound}
本节覆盖最终\texttt{field}、官方演练所用原\texttt{refined}及历史\texttt{enhanced}策略：问题三均用原点及半径1200的六点环，问题四前两者用25点双环、后者用990米31点三角格。联合调度只重排任务，已知频道筛选不增加补测义务；均保留每局部任务至多6次固定测向和原矩形光学后备。\texttt{field}每任务另允许至多一次原地预检测，问题三负观测的保守凸外包不扩大原域、保证清除的安全接近不增加当次接入距离。假定接口正常、误差模型成立且无现实时间中断。以下沿用较大的31点范数与次数界，仅为新预检测增加检测预算；不涵盖条带光学或自适应候选。此界不声称紧或最优，距离单位均为米。

\paragraph{可行域尺寸与216单元证书。}
记真源为$g$，最近正示向点为$s$，$\alpha=1.005^\circ$。示向轴坐标满足$0\le x\le1500$、$|y|\le x\tan\alpha$。计入显式外扩，令$\varepsilon=10^{-7}$，取
\[
 X=1500+2\varepsilon,\qquad
 W=2\bigl((1500+\varepsilon)\tan\alpha+\varepsilon/\cos\alpha\bigr).
\]
于是纵向长度至多$X$，横向宽度至多$W<53$，$\sqrt{X^2+W^2}<1501$。最小包围圆圆心$z\in P$，故$\operatorname{diam}(P)<1501$且$\|z-s\|<1501$。正测点范数至多3300、初始外切多边形范数小于1801，使相关\texttt{\_padding}为$10^{-9}<\varepsilon$；此为几何与显式外扩的解析界，不是全部浮点运算的形式化认证。

单点域直接用一圆覆盖。否则设直径端点差为$(a,b)$、$d=\sqrt{a^2+b^2}>0$。任意点差满足$|\Delta x|\le d$、$|\Delta y|\le W$、$|b|\le W$，故直径法向投影差$|-b\Delta x+a\Delta y|/d\le2W$。计入投影外扩，直径对齐矩形仍长小于1501、宽小于106，对角线小于1505。令$h_-=\sqrt2(19.9-8\varepsilon)$，不大于代码单元边长上限，则
\begin{equation}
 J\le\left\lceil\frac{1501}{h_-}\right\rceil
       \left\lceil\frac{106}{h_-}\right\rceil
   =54\times4=216
\end{equation}
个半径19.9的圆盘覆盖整个矩形。蛇形相邻中心距至多$\sqrt2\,19.9$，循环移位至多增加一条矩形内跳跃，故
\begin{equation}
 L_{\rm snake}\le(216-1)\sqrt2\,19.9+\sqrt{1501^2+106^2}<7570.
\end{equation}
每个覆盖中心$q_{\rm opt}$与$g$均在同一矩形中，故$\|q_{\rm opt}-g\|<1505$。

\paragraph{恢复点必须以最近正测点为中心估计。}
固定候选$q_0=z\pm\delta n$满足$n\perp(z-s)$、$\delta\le120$，因此$\|q_0-s\|\le A:=\sqrt{1501^2+120^2}<1506$。
在获得新正示向前，所有恢复仍使用同一个$s$。设$R_k$为相应轴向反射，则
\[
 q_{k+1}-s=\tfrac12R_k(q_k-s),\qquad
 \|q_{k+1}-s\|=\tfrac12\|q_k-s\|.
\]
反射只保持到$s$的距离，不能把失败点的源距与1500取平均。故连续恢复有$\|q_k-s\|\le2^{-k}A$；新正示向更新$s$后重用首点界。由$\|s-g\|\le1500$，全部局部测向点满足
\begin{equation}
 \|q-g\|\le B:=1500+A<3006,\qquad
 \|q\|\le1800+B<4806.
\end{equation}
圆心源距小于1501，覆盖中心源距小于1505，两者范数亦小于4806。测向间连线至多$2B<6020$；额外圆心清除接入至多$B+1501<4520$；光学接入至多$B+1505<4520$，若紧接圆心试清则至多$1501+1505<4520$。

\paragraph{任务进出、动作计数与总界。}
每任务至多6次固定测向，\texttt{field}在进入该过程前另可执行一次原地预检测；各提前返回分支均成功清除或转入有限光学覆盖。真源在覆盖域内，故每个正常完成任务必清除一新源，任务数至多16。成功终点距该源至多20、范数至多1820；原地预检测、共享补测及其\texttt{near}清除不移动。这处理了从\emph{另一源清除点}进入下一任务的情形。

骨架点范数至多$990\sqrt7<2620$。集成策略整次骨架扫描不插入局部移动，故至多31次骨架抵达，出发于原点、骨架或成功清除点，每段小于5240。任务首个需移动的动作由范数至多2620的任务外位置进入范数小于4806的局部点，入口小于7500；原地预检测不改变入口位置。此后至多5条测向连线、一次额外圆心清除移动（保证清除或小区域试清）、一次光学接入及一条蛇形，故
\begin{equation}
 L\le31\times5240+
 16\bigl(7500+5\times6020+4520+4520+7570\bigr).
 \label{eq:corrected-travel-bound}
\end{equation}
若首动作即圆心清除或光学首点，入口已支付该段，多余预算仅为放宽。任务至下一骨架的退出段已计入31次骨架抵达；无需返回原锚点。

骨架至多$31\times20$次测向；共享只在任务结束后调用，至多16次、每次至多16源；固定局部测向至多$16\times6$次。\texttt{field}原地预检测在六次之外，每任务至多一次，另加16次且不新增行程。放宽共享预算得到
\[
 M\le31\times20+16\times16+16\times6+16
 \le31\times20+31\times16+16\times6+16=1228.
\]
每源无信号后试清至多6次、小区域试清至多一次、覆盖至多216次，另留一次给非覆盖成功清除（骨架、共享、局部\texttt{near}或保证清除）；成功后不再处理，故$C\le16(216+6+1+1)=3584$。移动速度为5米/秒，测向含换频至多6秒，清除至多5秒，因而
\begin{equation}
 \begin{split}
 T&\le\frac{31\times5240+
 16(7500+5\times6020+4520+4520+7570)}5\\
 &\quad+1228\times6+3584\times5
 <100\times3600\ \mathrm{s}.
 \end{split}
\end{equation}
这是虚拟时间保守界，不是预期成绩或最优界；\texttt{refined}/\texttt{enhanced}没有新增预检测，仍可使用原$M\le1212$。\texttt{field}的任务重排及域缩小不改变其余范数、行程和光学次数预算，上界也不表示每个实例更快。现实20分钟限制须由客户端期限机制另行遵守。
